\section{Four Design Axes for Dynamics-Constrained Planning}

In line with the state of the art, my prior work reveals insights across four critical dimensions that govern robotic manipulation research. I also discuss the potential of learning-based approaches in addressing some of these challenges.

\header{Modeling Burden} While embodiment-driven analytical models provide theoretical foundations, they often fall short when modeling complex interactions. The progression from simulation-in-the-loop approaches (PokeRRT \cite{pasricha2022pokerrt}) to transformer-based dynamics models \cite{reed2024transformer} illustrates that learning-based methods have the potential to model non-linear and stochastic phenomena across diverse real-world settings. However, the effectiveness of these models remains highly task-dependent, as evidenced by the specialized constraints needed for liquid handling \cite{abderezaei2024clutter} versus whole-body manipulation under joint failures \cite{briscoe2024exploring}. Additionally, model fidelity and speed inform planning success across the high-dimensional search space of robot-object interactions in the real-world.

\header{Problem Dimensionality}
This search space involves the consideration of higher-order states like velocities, accelerations, and forces for dynamic motions, in addition to geometric constraints like positional limits and obstacle avoidance \cite{schmerling2019kinodynamic}. Finding good initial solutions that can then be optimized \cite{kuntz2019fast}, pre-computation in the combined robot and task space \cite{shome2021asymptotically,briscoe2024exploring}, and lazy approaches to dynamics propagation and collision checking \cite{pasricha2024virtues} can begin to address computational tractability issues especially as they arise from imperfect models in simple environments \cite{pasricha2022pokerrt,briscoe2024exploring,abderezaei2024clutter}.
However, even with recent advances in parallel collision-checking and forward kinematics calculations \cite{sundaralingam2023curobo,thomason2024motions}, computational complexity and memory consumption remain prohibitive for high-dimensional problems due to the inherent difficulty of generating valid samples, approximating distances in non-Euclidean spaces, and connecting states with dynamically feasible trajectories.
Learning-based approaches provide an alternative solution by implicitly compressing high-dimensional spaces through learned representations, which enables efficient and complex reasoning in state-action spaces \cite{qureshi2020motion}.

\header{System Design}
While traditional robotic planning systems offer verifiable safety assurance, explainable components for debugging, and generalizability to different scenarios, they face significant integration challenges and require extensive engineering effort when combining multiple specialized components \cite{abderezaei2024clutter}. Each component requires careful tuning, creating complex error propagation through multiple components, resulting in brittleness at system boundaries. Moreover, system performance typically does not improve without manual intervention and overall processing speed is subject to high variance depending on component choice and tuning.
Learning-based approaches offer advantages through end-to-end trainable architectures that naturally handle component interactions, improving scalability to new tasks without requiring extensive redesign as well as lending themselves to progressive improvement over the course of system deployment.

\header{Usability}
Traditional approaches require substantial engineering hours for each novel task and resulting edge cases---developing task-specific sampling strategies \cite{abderezaei2024clutter}, designing appropriate cost functions \cite{pasricha2024virtues}, and determining which physical phenomena to model in multi-objective scenarios \cite{pasricha2022pokerrt}. While traditional methods provide safety, completeness, and optimality guarantees, their implementation complexity often restricts deployment to controlled environments with simplified dynamics.
Learning-based methods significantly reduce this engineering burden by automatically extracting relevant patterns from demonstration or interaction data, where cost functions may be hard to specify.
However, high-quality training data is hard to acquire and limited by human demonstrations (typically geometric in nature), restricting the class of behaviors that can be learned by existing approaches.

These four dimensions highlight the opportunity learning-based methods present over kinodynamic planning for advancing robotic manipulation beyond geometric pick-and-place operations. Specifically, the integration of generative models with optimization-based approaches represents a promising direction for creating motion planners that efficiently reason about complex dynamics through pretrained representations of sensory input while generating trajectories that satisfy kinematic and dynamics constraints.